\documentclass[11pt]{article}
\usepackage[utf8]{inputenc}
\usepackage{amsmath,amssymb}
\usepackage[margin=1in]{geometry}
\usepackage{mathpazo}
\pagestyle{empty}
\begin{document}

June 9, 2026

Editor
Physical Review Letters

Dear Editor,

The standard theory of open quantum systems does not normally account for the topology of the causal graph connecting system and environment. Two systems with identical coupling Hamiltonians, identical bath spectral densities, and identical temperatures can differ in whether their interactions form closed causal loops -- and standard theory predicts identical behavior for both. We find the difference is not small.

A single causal cycle forces strictly positive quantum conditional mutual information for any $\theta \notin (\pi/2)\mathbb{Z}$, regardless of coupling strength. We prove that the QCMI vanishes if and only if every edge unitary in the cycle has its Cartan parameter in $(\pi/2)\mathbb{Z}$ -- that is, each gate is a Pauli rotation within the one-parameter family $\exp(i\theta\,\sigma\otimes\sigma)$. Any continuous deviation from these discrete values leaves the cycle unable to close, trapping quantum information regardless of coupling strength. The theorem is an "if and only if" -- an exact necessary and sufficient condition, placing it in a small class of exact results in open quantum dynamics.

The second result concerns the functional form: the correct scaling for aligned Cartan axes is $\theta^2\ln(1/\theta)$. A naive Taylor expansion of the von Neumann entropy would give $\theta^4$, missing the logarithmic singularity that dominates at small angles. The discrepancy exceeds a factor of 100 at accessible angles -- a factor of 19.4 at $\theta=\pi/8$ and 149 at $\theta=\pi/16$ -- verified by numerical Gram matrix diagonalization across six decades in the Cartan angle.

Two independent IBM Q measurements provide experimental support for the CFOL prediction. Protocol 1 (Z-basis, $2\times10^4$ shots): $\Delta I_Z>0$ at $\theta=\pi/4$ ($7\sigma$), $\pi/8$ ($4\sigma$), and $\pi/16$ ($2\sigma$, supporting evidence), with monotonic $\theta$-dependence. Protocol 2 (Y-basis symmetric ring, $10^5$ shots): $\Delta I_Y = 0.645$ ($129\sigma$), $0.354$ ($71\sigma$), and $0.128$ ($26\sigma$) bits at $\theta=\pi/4,\pi/8,\pi/16$. The experimental ratio $R_{\rm exp}=0.55$ exceeds the $\theta^4$ prediction of $0.063$ by an order of magnitude and matches $\theta^2\ln(1/\theta)$ scaling. The two protocols use different measurement bases and gate sequences, providing cross-validation against systematic errors. Statevector simulations predict two further protocols for the 4-node spatial ring (requiring native 4-cycle hardware): the ratio-test prediction gives $R=0.483$ vs $0.0625$ for naive $\theta^4$ scaling (7.7-fold gap), and $2.06\times$ QCMI amplification under axis misalignment. Complete circuit specifications appear in the Supplemental Material.

This manuscript establishes a rigorous theoretical connection between causal graph topology and quantum non-Markovianity, with experimental support from IBM Q hardware and a clear pathway to definitive verification on square-lattice processors.

We note two scope limitations: the theorem is proved for qubit systems (the Cartan subalgebra generalization to $d>2$ remains open), and it assumes a product initial environment state -- the physically natural setting for independent environmental degrees of freedom.

We confirm that this work is original, has not been published elsewhere, and is not under consideration by any other journal. The author has reviewed and approved the manuscript. The authors declare no competing interests.

\begin{flushright}
Respectfully submitted,\\[1.5em]
Zhongchang Huang\\
Independent Researcher\\
Nanjing, China
\end{flushright}

\end{document}
